\scanpage{140}
% archive foliation: 139
de $\unsure{\mathrm{Gl}|\mathrm{Sl}}(2,\mathbb Z)$ sur les structures de Torelli, par
\[
g\cdot u=u\circ g^{-1}\qquad
(u:\mathbb Z^2\xrightarrow{\sim}\pi_E).
\]

\uline{Corollaire 1} On a un iso. canonique
\begin{equation}
\mathcal T_{1,1}\simeq\mathrm{Sl}(2,\mathbb Z)
\tag{26}
\end{equation}
(où $\mathcal T_{11}=\pi_1(M_{1,1})$ est calculé relativement au revêtement universel choisi $\mathcal D$ de $M_{1,1}$).
\inserted{On a un iso. can
\begin{equation}
M_{1,1}\simeq(\unsure{\mathcal D},\mathrm{Sl}(2,\mathbb Z))
\tag{27}
\end{equation}}

Utilisant les grands théorèmes \unsure{marteaux-pilons}, on en déduit

\uline{Corollaire 2} Soit \deleted{$(E,a)$} une surface orientée \inserted{ponctuée} compacte connexe de genre $1$, et soit $\mathcal T$ le groupe des classes d'isotopie des automorphismes de $(E,a)$. \deleted{Alors considérons}\inserted{Alors} l'homomorphisme naturel
\begin{equation}
\mathcal T\longrightarrow\operatorname{Aut}_{\mathbb Z}^{+}(H)
\qquad\text{où }\quad H=\pi_1(E,a)
\tag{27}
\end{equation}
est un isomorphisme. $\unsure{\S}$

\uline{Corollaire 3} Soit $\Pi$ un groupe à lacets du type $\unsure{g}(1,1)$, $H=\Pi_{\mathrm{ab}}$ (qui est un $\mathbb Z$-module libre de rang $1$), $\mathcal T=\operatorname{Autext}^{+}_{\mathrm{lac}}(\Pi)$, \deleted{considérons}\inserted{alors, l'hom}
\begin{equation}
\mathcal T\longrightarrow\operatorname{Aut}_{\mathbb Z}^{+}(H)
\tag{28}
\end{equation}
est iso.

\margin{Il faudrait expliciter l'action \unsure{extérieure} de $\operatorname{Aut}_{\mathbb Z}^{+}(H)$ sur $\Pi$ via \unsure{l'action|l'identif.} $\operatorname{Aut}_{\mathbb Z}^{+}(H)\ldots$}
